% Generated draft: requires mathematical and source review.
\documentclass[11pt]{article}
\usepackage[margin=25mm]{geometry}
\usepackage{fontspec}
\setmainfont{Latin Modern Roman}
\usepackage{amsmath,amssymb,mathtools,unicode-math}
\setmathfont{Latin Modern Math}
\usepackage{enumitem,needspace,xurl,hyperref}
\setlength{\parindent}{0pt}
\setlength{\parskip}{5pt}
\setlength{\emergencystretch}{3em}
\newcommand{\N}{\mathbb{N}}
\newcommand{\Z}{\mathbb{Z}}
\newcommand{\Q}{\mathbb{Q}}
\newcommand{\R}{\mathbb{R}}
\newcommand{\C}{\mathbb{C}}
\newcommand{\F}{\mathbb{F}}
\newcommand{\A}{\mathbb{A}}
\newcommand{\PP}{\mathbb{P}}
\newcommand{\E}{\mathbb{E}}
\newcommand{\eps}{\varepsilon}
\newcommand{\dd}{\,\mathrm{d}}
\begin{document}
\section*{UnsolvedMath ID: 30000307}
\textbf{Source problem code:} catalogue entry
\subsection{Definitions and mathematical statement}
Retain \((a,b,c,d)\) for the parameters of the source's branching-coalescing (BRACO) process \(X\), and write \(Y\) for its corresponding resampling-selection process. This distinguishes the two processes whose symbols collided in the original extraction. Let \(P_t\) and \(Q_t\) denote their respective Markov semigroups in the intended state spaces. An invariant probability law \(\mu\) for \(X\) satisfies \(\mu P_t=\mu\) for every \(t\geq0\), and analogously for \(Y\).

The first question asks whether every invariant law of \(X\) is the Poissonization of an invariant law of \(Y\):
\[
 \forall\mu\ [\mu P_t=\mu\text{ for all }t\geq0]
 \ \Longrightarrow\ \exists\nu\ [\nu Q_t=\nu\text{ for all }t\geq0,
 \ \mu=\mathcal P\nu].
\]
Here \(\mathcal P\) is the specific Poissonization kernel in the report. In a usual countable-site construction, conditional on a state \(y\) sampled from \(\nu\), the particle counts at sites \(i\) are independent Poisson variables with intensities \(\lambda_i(y)\); the resulting mixture is \(\mathcal P\nu\). The intensity map, including any parameter-dependent normalization, is not supplied and is not set equal to \(y\) by assumption. The assertion ranges over all invariant laws in the intended state space, not merely ergodic, translation-invariant or moment-bounded laws unless the missing setup restricts that space. It is a surjectivity question and does not assert uniqueness of \(\nu\).

The second question concerns the reversed jump rates
\[
 a^\dagger(i,j)=a(j,i).
\]
Keeping \(b,c,d\) unchanged, does survival of the \((a,b,c,d)\) process imply survival of the \((a^\dagger,b,c,d)\) process? This is transposition of the prescribed rates, without additional stationary weights. If survival means global survival from a prescribed nonzero finite initial configuration \(x\), its usual event is \(\{\tau_0=\infty\}\), where \(\tau_0=\inf\{t\geq0:X_t=0\}\), with positive probability. The source does not specify that initial-state quantifier or distinguish global survival from local survival at a site; no one convention is silently substituted. Admissibility and nonexplosion of both original and reversed-rate processes require the missing hypotheses on the site set and rates.

The source does not provide either generator, the state spaces, rate summability or spatial assumptions, or the parameter relation defining the corresponding diffusion. These are essential for a completely specified theorem. The literature triage checked 2026-08-21 reports both invariant-law and survival questions explicitly posed and no general resolution located. The verification field reports repairing the collided process labels and defining the reversed rates. These are retained as supplied metadata without independent verification.
\subsection{Short English statement}
First, ask whether every stationary BRACO distribution comes from a stationary resampling-selection distribution by the source's Poisson mixture. Second, ask whether positive survival persists when migration rates are transposed. The generators, Poisson intensity normalization and survival convention are missing, while the two requests and the exact reversal formula are preserved.
\subsection{Sources}
\begin{itemize}
\item Local UnsolvedMath registry, permanent ID 30000307, source code catalogue entry, ``Invariant Laws and Survival of BRACO Processes,'' including original collided notation, repaired process identities and the two cleaned questions.
\item \emph{Mathematical Population Genetics (2006)}, Oberwolfach Reports, pp.\ 2241--2304, DOI 10.4171/owr/2005/40: \url{https://doi.org/10.4171/owr/2005/40}. This is the sole supplied literature reference, labelled ``original Oberwolfach report.'' The title uses 2006, while the DOI and proposed-year metadata use 2005.
\item Literature triage checked 2026-08-21; registry record updated 2026-09-16. No external source has been visited or independently reviewed here.
\end{itemize}
\end{document}
